\section{Coordinate variants and bounded paths}\label{sec:coordinate-variants}
\subsection{Projections and height conventions}

The coordinate arguments also apply to restrictions of the forest and
to other choices of finite component heights. We first record these
variants, then use bounded root paths to rule out an equivalent AUC norm.

\begin{lemma}\label{foundation:order-convex}
Let $E$ be a set of vertices whose intersection with every segment is
empty or a segment. The coordinate projection $P_E$ is contractive on
both $J$ and $X=J^*$. Its range on $J$ is the completion of $c_{00}(E)$
under the segment norm restricted to $E$, and
\[
 \|P_Ex\|=\|x|_{P_EJ}\| \qquad (x\in X).
\]
In particular these assertions apply to an ancestral set, its
complement, a difference of nested ancestral sets, a descendant subtree,
and an antichain. On an antichain the coordinate norm on either space
is exactly the usual $\ell_2$ norm.
\end{lemma}
\begin{proof}
For a finite vector $u$, each segment sum of $P_Eu$ is the sum of $u$
over the intersection of that segment with $E$. Nonempty intersections
are disjoint segments when the original segments are disjoint. Taking
the supremum proves $\rho(P_Eu)\le\rho(u)$. Extend by density and take
the adjoint to obtain the assertion on $X$. Contractivity and density
identify $P_EJ$ with the indicated completion. Since $P_E$ maps the
unit ball of $J$ onto the unit ball of its range, restriction has exactly
the asserted dual norm. A segment meets an antichain in at most one
vertex; singleton tests then give the $\ell_2$ norm on its primal range.
The restriction identity and $\ell_2$ duality give the dual assertion.
\end{proof}

\begin{proposition}\label{foundation:height-variants}
Let $I$ be countable and let $(h_i)_{i\in I}$ be nonnegative integers.
On the disjoint forest
\[
 \mathcal F_{\boldsymbol h}
   =\coprod_{i\in I}(\{i\}\times\mathbb N^{\le h_i})
\]
define $J_{\boldsymbol h}$ by the same disjoint-segment norm and set
$X_{\boldsymbol h}=J_{\boldsymbol h}^*$. If $J(h)$ denotes the
single-tree completion on $\mathbb N^{\le h}$, then
\[
 J_{\boldsymbol h}
   =\Big(\bigoplus_{i\in I}J(h_i)\Big)_{\ell_2},
 \qquad
 X_{\boldsymbol h}
   =\Big(\bigoplus_{i\in I}J(h_i)^*\Big)_{\ell_2}
\]
isometrically under coordinate restriction. Both spaces are reflexive,
and finite coordinate vectors are dense in each. The projection and
complementary-subtree formulas remain valid for this forest.

The convention of one tree with $n$ levels, including the root, for
every $n\ge1$ gives exactly
\[
 J_{\mathrm{levels}}=\mathbb R\oplus_2 J,
 \qquad X_{\mathrm{levels}}=\mathbb R\oplus_2 X.
\]
Its component estimate is
$\|u\|_2\le\rho(u)\le\sqrt n\,\|u\|_2$, and a root-to-leaf path
has $n-1$ extension steps.
\end{proposition}
\begin{proof}
Every segment belongs to one component. Grouping a test family by
component gives one inequality in the primal sum formula. For the
other, choose nearly maximizing finite families on each of the
finitely many components meeting the support and combine them.
Completion gives the displayed primal identity, and duality for
$\ell_2$ sums gives the dual identity. On component $i$,
\[
 \|u\|_2\le\rho(u)\le\sqrt{h_i+1}\,\|u\|_2.
\]
Thus each component and its dual are reflexive Hilbert space
renormings with dense coordinates. The sum identities give
reflexivity and coordinate density on the entire forest. No bound
uniform in $i$ on the Hilbert space comparison is needed. Projection
contractivity uses only segment intersection, as in
Lemma~\ref{foundation:order-convex}; the complementary-subtree identity
uses independent segment tests in its incomparable subtrees.

For the last assertion, put $h=n-1$. The component $n=1$ is the
one-dimensional singleton tree with the absolute-value norm. The
remaining components have heights $1,2,\ldots$, and the coordinate map
$(n,\sigma)\mapsto(n-1,\sigma)$ preserves their segments and all their
tests. Splitting off the singleton proves both isometries. The
component bound and the number of extension steps follow from the
maximal chain cardinality $n$.
\end{proof}

\subsection{Bounded paths and equivalent norms}

A single segment functional has norm one even when the segment is very
long. Its successive coordinate increments therefore give bounded
paths. The following argument explains why weakly null choices of the
next increment obstruct an asymptotically uniformly convex equivalent
norm. It also retains the precise radius when increments are normalized
before being approximated in a finite-codimensional subspace.

For a norm $N$ on an infinite-dimensional real Banach space $Z$, write
$S_N=\{z:N(z)=1\}$ and let $\operatorname{cof}(Z)$ denote the closed
finite-codimensional linear subspaces. Define
\begin{equation}\label{eq:one-sided-modulus}
 \overline\delta_N(t)=\inf_{x\in S_N}
   \sup_{F\in\operatorname{cof}(Z)}
   \inf_{y\in F\cap S_N}\bigl(N(x+ty)-1\bigr),\qquad t>0.
\end{equation}
Positivity at every $t>0$ is asymptotic uniform convexity.
The spheres $F\cap S_N$ are nonempty. A norming functional for $x$,
with $F$ contained in its kernel, gives the lower bound in
$0\le\overline\delta_N(t)\le t$; the upper bound is the triangle
inequality.

\begin{lemma}\label{foundation:weak-null-growth}
Let $x\in S_N$, let $(z_j)$ be weakly null, and suppose $N(z_j)\ge t>0$
for every $j$. For every
$0<\gamma<\overline\delta_N(t)$, all sufficiently large $j$ satisfy
\[
 N(x+z_j)>1+\gamma.
\]
\end{lemma}
\begin{proof}
Choose $\eta$ with $\gamma<\eta<\overline\delta_N(t)$ and then
$F\in\operatorname{cof}(Z)$ such that
$N(x+tv)>1+\eta$ for every $v\in F\cap S_N$.
Put $u_j=z_j/N(z_j)$. These vectors are weakly null because
$1/N(z_j)\le1/t$. Their images in $Z/F$ converge to zero weakly,
hence in norm, since the quotient is finite dimensional. Choose
$v_j\in F$ with $N(v_j-u_j)\to0$. Then $N(v_j)\to1$, so
$w_j=v_j/N(v_j)$ is defined for large $j$, belongs to $F\cap S_N$,
and satisfies $N(w_j-u_j)\to0$. Consequently
\[
 N(x+t u_j)>1+\eta-tN(w_j-u_j)>1+\gamma
\]
for large $j$. For $r\ge t$, convexity gives
\[
 N(x+r u_j)-1\ge\frac rt\bigl(N(x+t u_j)-1\bigr).
\]
Apply this at $r=N(z_j)$ to finish the proof.
\end{proof}

There is a second useful form of the approximation step. If
$N(z_j)\ge2t$ and $N(x+tv)>1+\eta$ on $F\cap S_N$, approximate
$z_j$ directly by $v_j\in F$ with $N(z_j-v_j)\to0$.
Eventually $N(v_j)\ge t$, so the same ray-convexity inequality gives
$N(x+v_j)>1+\eta$. For every $0<\gamma<\eta$, choosing
$N(z_j-v_j)<\eta-\gamma$ gives $N(x+z_j)>1+\gamma$.
This form uses the spare factor two in the radius and does not require
approximating the normalized sequence.

\begin{proposition}\label{foundation:bounded-tree}
Suppose an infinite-dimensional real Banach space $(Z,\|\cdot\|)$
contains, for every positive
integer $k$, vectors $(u_s)_{s\in\mathbb N^{\le k}}$ such that
\[
 0<m\le\|u_s\|\le M,
 \qquad \|u_{s^\frown j}-u_s\|\ge c>0,
\]
where $m,M,c$ are independent of $k$, and for each nonterminal $s$ the
sequence $(u_{s^\frown j}-u_s)_j$ is weakly null against the whole
continuous dual of $Z$. If
$0<a\le b<\infty$ and $a\|z\|\le N(z)\le b\|z\|$, then
\[
 \overline\delta_N\!\left(\frac{ac}{bM}\right)=0.
\]
In particular the same conclusion holds at $ac/(2bM)$, and $Z$ admits
no equivalent asymptotically uniformly convex norm.
\end{proposition}
\begin{proof}
Suppose the modulus at $t=ac/(bM)$ is positive and choose
$0<\gamma<\overline\delta_N(t)$. At a nonterminal $s$, put
\[
 x=\frac{u_s}{N(u_s)},\qquad
 z_j=\frac{u_{s^\frown j}-u_s}{N(u_s)}.
\]
The equivalent norms have the same continuous linear functionals,
so $(z_j)$ is weakly null, and $N(z_j)\ge ac/(bM)=t$.
Lemma~\ref{foundation:weak-null-growth} gives a child for which
$N(u_{s^\frown j})>(1+\gamma)N(u_s)$. Starting at the root and
repeating $k$ times gives $am(1+\gamma)^k<bM$, which is impossible
for sufficiently large $k$. Thus the modulus at $t$ is zero. Applying
the same argument at $t/2$, using either the lemma or the direct
approximation form above, proves the stated smaller-radius conclusion.
\end{proof}

\begin{corollary}\label{foundation:forest-obstruction}
If the finite component heights $h_i$ in
Proposition~\ref{foundation:height-variants} are unbounded, then every
equivalent norm $N$ on $X_{\boldsymbol h}$ with
$a\|x\|\le N(x)\le b\|x\|$ satisfies
\[
 \overline\delta_N(a/b)=\overline\delta_N(a/(2b))=0.
\]
\end{corollary}
\begin{proof}
For a vertex $s$ in a component with root $o$, put
$u_s=\sum_{o\preceq v\preceq s}e_v^*$. A single segment test gives
$\|u_s\|\le1$, while evaluation on $e_s$ gives equality. Likewise
$\|e_s^*\|=1$, and
$u_{s^\frown j}-u_s=e_{s^\frown j}^*$.
These increments are weakly null: under the canonical identification
$X_{\boldsymbol h}^*=J_{\boldsymbol h}$, their evaluations on $g$
are the coordinates $g(s^\frown j)$, which tend to zero since
$\|g\|_2\le\rho(g)$. Equivalently, the coordinate map on this fixed
component is a Hilbert space isomorphism, so every continuous
functional on the full space restricts to a continuous Hilbert space
functional there. Components of unbounded finite height give the
trees required by Proposition~\ref{foundation:bounded-tree}, with
$m=M=c=1$. For the convention of $n$ levels the same iteration makes
$n-1$ steps. Both step counts tend to infinity.
\end{proof}
